\section{Technique}
\label{sec:technique}

This section first gives the background the paper needs
(Section~\ref{sec:background}) and defines the problem
(Section~\ref{sec:problem}). It then describes the workflow of \tool and its
four steps.

\section{Background}
\label{sec:background}

\paragraph{Turso.} Turso~\cite{turso} reimplements SQLite in Rust and aims to
accept the same SQL. Its shell, \texttt{tursodb}, reads statements on standard
input. Text and blob results are bounded by a constant,
\texttt{Value::MAX\_BLOB\_LENGTH}. A statement whose result would exceed it should
fail with the error \emph{string or blob too big}, which is SQLite error code~18.
Rust aborts the process when an allocation fails and cannot be recovered, so a
size check that runs after the allocation is not a check at all.

\paragraph{Differential testing.} A \emph{differential oracle} runs one query on
two engines and reports a disagreement. It needs no knowledge of the correct
answer, but it needs a \emph{reference engine} that implements the same feature.
Turso's fuzzer uses SQLite as its reference engine. SQLite has no \sql{lpad},
\sql{rpad} or \sql{repeat}, so a disagreement on those functions cannot be
judged.

\paragraph{Crash oracle.} A \emph{crash oracle} reports a statement that ends the
engine's process, panics, reports an internal error or does not answer within a
time limit. It does not report a \emph{user error}, which is a statement the
engine rejects with a normal message, such as a syntax error or a documented
limit. A crash oracle needs no reference engine. It finds no wrong answer that
leaves the process running.

\paragraph{Bug reports as test material.} DuckDB maintainers add the label
\emph{reproduced} to an issue once they have confirmed it. A fix usually adds a
regression test that contains the query which triggered the issue. Such a query
is small, it was chosen by someone who understood the defect, and it tends to sit
in an awkward corner: an empty string, a multi-byte character, a value at the
edge of a type.

\paragraph{Literal and template.} We call a constant in a query, such as
\sql{10} or \sql{'x'}, a \emph{literal}. A \emph{template} is a query in which
every literal is replaced by a typed hole, such as \sql{\{n0\}} for an integer or
\sql{\{s4\}} for a string. Filling the holes gives a new query with the same
structure.

\section{Problem Statement}
\label{sec:problem}

We now state the problem precisely. Let \Fz be the fuzzer that already tests an
engine~$E$, and let \Rengine be the reference engine that \Fz compares against. We
write \Voc for the set of SQL functions that \Fz can emit, and \Lit for the
literals it can emit: a range of integers and a set of strings. We write \VocR
for the functions that \Rengine implements.

A statement $q$ is \emph{expressible} by \Fz if every function in $q$ is in \Voc
and every literal in $q$ is in \Lit. A statement is \emph{judgeable} if every
function in $q$ is in \VocR. The \emph{blind spot} of \Fz is
\begin{equation}
  \label{eq:blind}
  \Blind = \{\, q \mid q \text{ is not expressible by } \Fz
           \ \lor\ q \text{ is not judgeable against } \Rengine \,\}.
\end{equation}
\Fz never generates a statement outside its own vocabulary or literal domain,
however long it runs. A statement that \Fz does generate but cannot judge is
generated in vain. Both halves of~\eqref{eq:blind} are properties of the source
code of \Fz and \Rengine, so \Blind can be measured before any test runs.

\paragraph{Goal.} Find statements $q \in \Blind$ on which $E$ violates the crash
oracle, without a reference engine, and with each finding regenerated from a
recorded seed. We look for two things: which functions and values put a
statement in \Blind, and where to get statements in that region that are likely
to break $E$.

\paragraph{Why not widen the incumbent fuzzer instead?} Its maintainers could add
\sql{lpad} to \Voc and \imax to \Lit, and the defects in this paper would then be
within reach. But that needs someone to know which function and which value to
add. The blind spot is large: any function outside \Voc and any value outside
\Lit belongs to it. Bug reports from another engine point at the corners where
someone has already been hurt. A crash oracle then removes the need for a
reference engine, so a function SQLite lacks can still be tested.

\paragraph{Why not run the other engine's tests directly?} A regression test was
written for the other engine's defect. Its literal values are tuned to that
defect, which the engine under test probably never had, and its syntax may not
parse. The structure of the query, meaning which functions are nested how, is
what carries over. The values have to be chosen again.


\subsection{Workflow}
\label{sec:workflow}

Figure~\ref{fig:overview} shows the workflow of \tool. The main idea is to aim
testing at the blind spot of the incumbent fuzzer, and to take the statements
that reach it from bug reports another engine's maintainers have already
reduced. Two kinds of input go in: the incumbent fuzzer's source, and a bug
tracker. One kind comes out: a statement that ends the engine's process, with the
seed that regenerates it. Step~3 loops over seeds and is the only expensive step.

\begin{figure}[t]
  \centering
  %% Figure: overview of GRAFT. Steps left to right; the loop over seeds is dashed.
%% Grayscale-safe (black lines, black!6 fill), \scriptsize text, step numbers
%% match the "Step N" subsections of the Approach.
\begin{tikzpicture}[
    x=1cm, y=1cm, font=\scriptsize,
    box/.style={draw, rounded corners=2pt, align=center, inner sep=2pt,
                minimum height=10mm, text width=21mm, fill=white},
    src/.style={draw, align=center, inner sep=2pt, fill=black!6,
                minimum height=8mm, text width=21mm},
    outbox/.style={draw, align=left, inner sep=3pt, fill=black!6,
                text width=27mm, font=\tiny},
    flow/.style={-{Stealth[length=1.6mm]}, line width=0.5pt},
    hint/.style={-{Stealth[length=1.6mm]}, line width=0.5pt, dash pattern=on 1pt off 1pt},
    back/.style={-{Stealth[length=1.6mm]}, line width=0.5pt, dash pattern=on 2pt off 1.5pt},
  ]
  % inputs
  \node[src] (fz)  at (0, 1.7)  {incumbent fuzzer\\source code};
  \node[src] (trk) at (0, 0)    {another engine's\\bug tracker};
  \node[src] (gen) at (3.35,-1.75) {typed SQL\\generator};
  % steps
  \node[box] (s1) at (3.35, 1.7) {\textbf{1} measure the\\blind spot};
  \node[box] (s2) at (3.35, 0)   {\textbf{2} mine reports,\\make templates};
  \node[box] (s3) at (6.7, 0)    {\textbf{3} fill with boundary\\values and run};
  \node[box] (s4) at (10.05, 0)  {\textbf{4} reduce, re-check,\\report};
  % outputs
  \node[outbox] (o1) at (6.7, 1.7) {functions and values the\\fuzzer cannot emit or judge};
  \node[outbox] (o4) at (10.05,-1.75) {defect: one statement,\\seed, signature, report};
  % edges
  \draw[flow] (fz) -- (s1);
  \draw[flow] (s1) -- (o1);
  \draw[flow] (trk) -- (s2);
  \draw[flow] (s2) -- (s3);
  \draw[flow] (gen.east) -| ([xshift=-4mm]s3.south);
  \draw[hint] (o1) -- node[right, font=\tiny] {where to aim} (s3);
  \draw[flow] (s3) -- node[above, font=\tiny] {crash} (s4);
  \draw[flow] (s4) -- (o4);
  \draw[back] ([xshift=3mm]s3.south) to[out=-80, in=-100, looseness=5]
        node[below, font=\tiny] {next seed} ([xshift=9mm]s3.south);
\end{tikzpicture}

  \caption{Workflow of \tool. Step~1 runs once and tells Step~3 where to aim.
  Steps~2 and~3 supply two sources of statements, mined templates and a typed
  generator. Step~3 loops over seeds and is the only step whose cost grows with
  the budget; a crash goes to Step~4.}
  \label{fig:overview}
\end{figure}

\subsection{Step 1: Measure the Blind Spot}
\label{sec:step-blind}
We read the incumbent fuzzer's source and extract three facts: (i)~its function
vocabulary \Voc; (ii)~its literal domain \Lit; and (iii)~its reference engine
\Rengine. For Turso these are a table of 47 scalar functions and a default literal
policy in two source files~\cite{tursofuzzer}. The policy sets
\texttt{int\_max} to 1{,}000{,}000 and restricts strings to alphanumeric
characters of at most 50 bytes. Every function outside \Voc, every integer outside
the range, every non-alphanumeric string, and every function SQLite lacks is in
\Blind by Equation~\eqref{eq:blind}. This step takes reading, not running. It
needs no new tool, and we do not claim reading source code as new. What we
claim is that this reading should come first.

\subsection{Step 2: Mine Reports and Make Templates}
\label{sec:step-mine}
\tool collects closed issues that carry the tracker's label for a confirmed
defect; for DuckDB this label is \emph{reproduced}. For each issue it fetches the
fixing commit and the regression test that the fix added, and extracts every SQL
statement from that test and from the issue body. Setup statements such as
\sql{CREATE TABLE} and \sql{INSERT} are kept with their query, so a reproduction
brings its own schema and data.

Each extracted query becomes a template. \tool replaces every literal with a
hole typed by the literal's kind (integer, real or string) and leaves everything
else exactly as written: function names, nesting, clauses and setup.
Figure~\ref{fig:motivating}b shows the result for one query. We keep the
structure fixed on purpose. Re-pointing a mined query at a new schema would need
column types and arities to match, and a wrong substitution produces SQL that
merely looks plausible. A query whose structure is unchanged is valid by
construction, so no repair step is needed. The cost is that \tool cannot explore
a structure that no report contains.

\subsection{Step 3: Fill with Boundary Values and Run}
\label{sec:step-run}
\tool fills each hole from a pool weighted toward the edges of each type. For
integers the pool holds zero, $\pm 1$, the limits of 8-, 16- and 32-bit integers,
and the limits of 64-bit integers, including \imax. String holes are filled from
a pool of short ASCII strings. Every random choice comes from a generator seeded
per case, and the seed is stored with each result.

Mined templates are one source of statements. The second is a typed SQL generator
that builds a schema, inserts rows, and composes expressions and queries. It
emits only features that the target engine accepted when \tool probed it, so it
rarely wastes a case on syntax the engine rejects. Its string pool also includes
multi-byte values such as \sql{'é'} and three-byte Chinese characters. A command
line flag selects either source or both.

Each case, meaning its setup statements and its query, runs in a fresh
\texttt{tursodb} process under a time limit and a memory limit. The crash oracle
classifies the outcome into one of six classes (\emph{none}, \emph{user error},
\emph{internal error}, \emph{panic}, \emph{crash}, \emph{timeout}) and reports
only the last four. We do not claim the crash oracle as new; fuzzers have used it
for decades. We chose it because it is the one oracle that can judge a function
the reference engine lacks.

\subsection{Step 4: Reduce, Re-check and Report}
\label{sec:step-triage}
A violation is not yet a finding. \tool groups violations by a \emph{signature},
a hash of the oracle, the error class and the error text with numbers and paths
removed. It then (i)~reduces the statement while keeping the same violation;
(ii)~re-runs it on a fresh process and on a build of the current upstream main
branch, so that a defect fixed last week is not reported; (iii)~searches the
upstream tracker and records the queries it used; and (iv)~drafts a report. A
person reads the draft, reproduces it, and decides whether to file it. Reduction
by repeated deletion is standard and we do not claim it as new. The duplicate
check matters more than it looks. In our campaign it withheld one candidate,
because an upstream pull request had already changed the code it touched.
